\begin{table}[htbp]
\centering
\caption{Main results on MVTec AD (5-task, 15 categories). I-AUROC (\%) after all tasks. MemoryAD uses DINOv2 ViT-B/14 with budget $B=10\text{K}$.}
\label{tab:main_results}
\begin{tabular}{lccccc}
\toprule
Method & I-AUROC $\uparrow$ & P-AUROC $\uparrow$ & FR $\downarrow$ & Avg Inc. & FT $\uparrow$ \\
\midrule
Joint (reference baseline) & 95.1 & -- & -- & -- & -- \\
\textbf{MemoryAD (ours)} & 94.0 $\pm$ 0.8 & 96.0 & 0.97 & 94.0 & 99.0 \\
UCAD (Liu et al. 2024)* & 93.0 & -- & -- & -- & -- \\
ONER (Li et al. 2024)* & 93.4 & -- & -- & -- & -- \\
Replay + RD4AD & 83.7 & -- & -- & -- & -- \\
Naive (lower bound) & 78.4 & -- & -- & -- & -- \\
EWC + RD4AD & 74.1 & -- & -- & -- & -- \\
LwF + RD4AD & 70.0 & -- & -- & -- & -- \\
\bottomrule
\end{tabular}
\vspace{1mm}
\footnotesize{Forgetting metrics (FR, Avg Inc., FT) for baselines are denoted as '--' because the original papers did not report the full $T \times T$ incremental evaluation matrix, and replicating their training-heavy CIL pipelines across all task boundaries is computationally prohibitive. Baselines were evaluated on final I-AUROC only. *Note: UCAD and ONER results are cited from their 15-class class-incremental literature values where available (evaluated on MVTec AD using a 15-task protocol with 1 category per task).}
\end{table}